\documentclass[11pt,a4paper]{article}

% ----------------------------
% Packages
% ----------------------------
\usepackage[utf8]{inputenc}
\usepackage[T1]{fontenc}
\usepackage{lmodern}
\usepackage{geometry}
\geometry{margin=1in}

\usepackage{amsmath,amssymb}
\usepackage{graphicx}
\usepackage{booktabs}
\usepackage{hyperref}
\usepackage{cite}
\usepackage{enumitem}

\setlist{nosep}

% ----------------------------
% Metadata
% ----------------------------
\title{Satellite Planning and Dimensioning: Project Report\\
\large Proofs of Concept, Market Exploration, and Lessons Learned}

\author{
Author Name \\
Institution / École d’Ingénieurs \\
Project / Grant Reference
}

\author{
Author Name \\
Institution / École d’Ingénieurs \\
Project / Grant Reference
}


\date{\today}

% ----------------------------
% Document
% ----------------------------
\begin{document}

\maketitle

\begin{abstract}
This report summarizes the objectives, technical work, market exploration, and outcomes of the project funded by BPI. 
We describe the initial hypotheses, the approaches investigated, experimental results (including negative findings),
and the implications for future development and industrialization.
\end{abstract}

\section{Executive Summary}
% This section provides a concise, non-marketing summary of:
% (i) the problem addressed,
% (ii) what was attempted,
% (iii) what was learned,
% (iv) what remains unresolved.
% It should be understandable without technical background.

\section{Problem Statement and Motivation}
\subsection{Technical Problem}
% Describe the core technical challenge.
% State clearly what was unknown or risky at project start.

\subsection{Market Opportunity}
% Describe the target context, users, and use cases.
% Explain why this problem matters economically or operationally.

\section{State of the Art}
\subsection{Existing Technical Approaches}
% Summarize relevant academic and industrial solutions.
% Clearly distinguish prior art from project contributions.

\subsection{Limitations of Existing Solutions}
% Explain why existing approaches are insufficient.
% This justifies the project’s necessity.

\section{Initial Hypotheses and Objectives}
% List the main technical and market hypotheses tested.
% Clarify success criteria and assumptions.

\section{Technical Approach}
\subsection{System Overview}
% Provide a high-level description of the architecture or methodology.

\subsection{Formal Model or Algorithm}
% When relevant, introduce mathematical notation.
% For example, let $x \in \mathbb{R}^n$ denote the input vector,
% $\theta$ the model parameters, and $f_\theta(x)$ the system output.
% Define all symbols explicitly.

\subsection{Implementation Details}
% Languages, tools, constraints, and design decisions.

\section{Product Specification: LEO Dimensioning Solution (v0.1)}
\subsection{Purpose}
Provide a dimensioning decision for a LEO satellite operator: the minimum number of satellites required to meet
operator-defined quality-of-service (QoS) targets, under priors on user demand/mobility and
satellite distributions. An example user could be an operator providing satellite internet.

\subsection{Users}
Day-to-day users are LEO satellite operators who own and operate satellite constellations.

\subsection{Decision and Optimisation Target}
The product outputs a recommended constellation size $N$ (number of satellites) such that QoS is satisfied with high
probability.
In v0.1, the optimisation objective is QoS feasibility (not cost, routing, or detailed PHY performance).

\subsection{Geographic Usage Zones (v0.1)}
In v0.1, the operator specifies a single latitude band:
\[
[\mathrm{lat}_{\min}, \mathrm{lat}_{\max}], \quad \mathrm{lat}_{\min} \le \mathrm{lat}_{\max}.
\]
Users are drawn only from within this band.

\subsection{Inputs (v0.1)}
\begin{itemize}
    \item Latitude band $[\mathrm{lat}_{\min}, \mathrm{lat}_{\max}]$.
    \item Priors for user distribution and mobility (e.g., moving users such as trains and planes).
    \item Priors for satellite distribution/constellation model (parameterised by $N$).
    \item QoS targets: minimum coverage and minimum throughput (see below).
    \item Certification parameters: confidence level $1-\alpha$ and feasibility target $p_{\mathrm{target}}$.
\end{itemize}

\subsection{Outputs (v0.1)}
\begin{itemize}
    \item Recommended minimum constellation size $N$.
    \item Certificate: a confidence-backed lower bound on $\mathbb{P}(\text{QoS met})$ under the stated priors.
    \item Reproducibility artifact: simulator configuration and seeds sufficient to replay the certification run.
\end{itemize}

\subsection{QoS Definition (v0.1)}
QoS is defined by two minimum constraints that must both be met:
\begin{itemize}
    \item \textbf{Coverage}: at least $C_{\min}$ fraction of users in the latitude band are served (exact definition of
    ``served'' to be specified by the simulator model).
    \item \textbf{Throughput}: at least $R_{\min}$ throughput is delivered (exact definition to be specified; in v0.1 this is
    expected to be driven by a simplified satellite physical resource block (PRB) capacity model).
\end{itemize}

\subsection{Statistical Certification (v0.1)}
Define a Bernoulli per-trial success event $A$ where $A=1$ iff both QoS constraints (coverage and throughput) are met in
that trial.
Let $p := \mathbb{P}(A=1)$ under the stated priors and simulator configuration.
The product certifies feasibility by returning, with confidence $1-\alpha$, a lower confidence bound $\mathrm{LCB}(p)$ and
declaring the design feasible if $\mathrm{LCB}(p) \ge p_{\mathrm{target}}$.

\subsection{Scope and Non-Goals (v0.1)}
\begin{itemize}
    \item \textbf{In scope}: QoS-driven constellation size dimensioning with certification; simplified capacity model via a
    PRB grid on satellites; minimal corresponding user properties needed to consume PRBs.
    \item \textbf{Out of scope}: detailed PHY/link-budget simulation; full network routing; end-to-end application performance.
\end{itemize}

\subsection{Implementation Constraint (v0.1)}
The v0.1 specification must be achievable with minimal changes to the existing codebase; planned additions are limited to
introducing a PRB grid on satellites (previously modelled as capacity-free points) and minimal corresponding user properties.

\section{Experiments and Results}
\subsection{Experimental Setup}
% Datasets, evaluation protocol, baselines, and metrics.

\subsection{Results}
% Present quantitative and qualitative results.
% Include tables and figures where appropriate.

\subsection{Negative Results and Limitations}
% Explicitly document what did not work and why.
% This section is important for funding evaluation.

\section{Market Validation Efforts}
% Describe interviews, pilots, or tests conducted.
% State what evidence was obtained and what remains inconclusive.

\section{Key Learnings}
% Summarize technical and market insights gained.
% Distinguish validated assumptions from rejected ones.

\section{Remaining Risks and Open Questions}
% Identify unresolved technical, operational, or market risks.

\section{Next Steps and Industrialization Path}
% Describe realistic next steps.
% Explain how additional funding or partnerships would reduce remaining risks.

\section*{Acknowledgments}
% Optional. Acknowledge funding support and collaborators.

\bibliographystyle{plain}
\bibliography{references}

\end{document}
